% answers_short.tex —— 2014 年·上海卷·理科 · 简短答案速查
% 由 tools/make_short_answers.py 自 answers.tex 的 \gksolhead/\gkanswer 提取，
% 并与 2014/2014_数学真题_上海理_答案要点.md 逐题比对后生成，勿手改（改源 MD 后重跑）。

\begin{gknotice}
\textbf{使用说明}：本节只给最终答案便于快速核对，完整推导见解析版。答案为编写者独立推导，\textbf{非官方标准答案}。
\end{gknotice}

\gkpart{填空题}
\begin{itemize}[leftmargin=2.6em]
  \item[\textbf{1.}] $\dfrac\pi2$。由 $y=1-2\cos^2(2x)=-\cos4x$ 得 $T=\dfrac{2\pi}4=\dfrac\pi2$。
  \item[\textbf{2.}] $6$。$\left(z+\dfrac1{\overline z}\right)\overline z=|z|^2+1=5+1=6$。
  \item[\textbf{3.}] $x=-2$。椭圆右焦点为 $(2,0)$，故 $\dfrac p2=2$，$p=4$，准线 $x=-\dfrac p2=-2$。
  \item[\textbf{4.}] $(-\infty,2]$。$a\leqslant2$ 时 $2\in[a,+\infty)$，$f(2)=2^2=4$；$a>2$ 时 $f(2)=2\neq4$。
  \item[\textbf{5.}] $2\sqrt2$。$x^2+2y^2\geqslant2\sqrt2|xy|=2\sqrt2$，$x^2=2y^2$ 且 $xy=1$ 时取等号。
  \item[\textbf{6.}] $\arccos\dfrac13$。由 $\pi rl=3\pi r^2$ 得 $l=3r$，故 $\cos\theta=\dfrac rl=\dfrac13$。
  \item[\textbf{7.}] $\dfrac13$。曲线即 $3x-4y=1$，令 $y=0$ 得 $x=\dfrac13$。
  \item[\textbf{8.}] $\dfrac{\sqrt5-1}{2}$。由 $a_1=\dfrac{a_1q^2}{1-q}$ 得 $q^2+q-1=0$，结合 $|q|<1$ 取正根。
  \item[\textbf{9.}] $(0,1)$。$x>0$，$x^{\frac23}<x^{-\frac12}\iff x^{\frac76}<1\iff0<x<1$。
  \item[\textbf{10.}] $\dfrac1{15}$。连续三天共 $8$ 种，$\dfrac8{\mathrm C_{10}^3}=\dfrac8{120}=\dfrac1{15}$。
  \item[\textbf{11.}] $-1$。必有 $a^2=b$，$b^2=a$，得 $a^2+a+1=0$，故 $a+b=a+a^2=-1$。
  \item[\textbf{12.}] $\dfrac{7\pi}3$。$2\sin\left(x+\dfrac\pi3\right)=a$ 恰三解当且仅当 $a=\sqrt3$，三解为 $0$，$\dfrac\pi3$，$2\pi$。
  \item[\textbf{13.}] $\dfrac15$。$4.2=E(\xi)\leqslant4+p_5$，故 $p_5\geqslant0.2$，且 $p_5=0.2$，$p_4=0.8$ 时可达。
  \item[\textbf{14.}] $[2,3]$。$A$ 为 $PQ$ 中点，$m=\dfrac{x_P+6}2$，而 $x_P\in[-2,0]$。
\end{itemize}

\gkpart{选择题}
\begin{center}\begin{tabular}{ccccc}
\toprule
题号 & 15 & 16 & 17 & 18 \\
\midrule
答案 & B & A & B & D \\
\bottomrule
\end{tabular}\end{center}

\gkpart{解答题}
\begin{itemize}[leftmargin=2.6em]
  \item[\textbf{19.}] $P_1P_2=P_2P_3=P_3P_1=4$，$V=\dfrac{2\sqrt2}3$。展开图为三角形迫使侧面为正三角形，侧棱 $l=2$，棱长全为 $2$。
  \item[\textbf{20.}] (1) $f^{-1}(x)=\log_2\dfrac{4(x+1)}{x-1}$，定义域 $(-\infty,-1)\cup(1,+\infty)$。由 $2^x=\dfrac{4(y+1)}{y-1}$ 解出；(2) $a=0$ 时为偶函数；$a=1$ 时为奇函数；$a>0$ 且 $a\neq1$ 时既非奇函数也非偶函数。$a>0$ 且 $a\neq1$ 时定义域 $\{x\mid x\neq\log_2a\}$ 不关于原点对称
  \item[\textbf{21.}] (1) 至多 $20\sqrt2$ 米，约 $28.28$ 米。由 $\tan\alpha\geqslant\tan2\beta$ 得 $6400-h^2\geqslant5600$；(2) 约 $26.93$ 米。由 $x=\dfrac{115\tan\beta}{\tan\alpha+\tan\beta}$、$h=x\tan\alpha$ 得水平偏差约 $0.69$ 米、高约 $26.92$ 米，$CD=\sqrt{0.69^2+26.92^2}$
  \item[\textbf{22.}] (1) 证明。$\eta=(1+2-1)(-1+0-1)=2\times(-2)=-4<0$，满足分隔定义；(2) $k\in\left(-\infty,-\dfrac12\right]\cup\left[\dfrac12,+\infty\right)$。$|k|<\dfrac12$ 时直线与双曲线相交；$|k|\geqslant\dfrac12$ 时无公共点，且取关于 $y$ 轴对称的两点可使 $\eta=u^2-k^2t^2<0$；(3) 证明。$y$ 轴与 $E$ 无公共点且分隔 $(1,2)$ 与 $(-1,2)$；其余过原点的直线 $y=kx$ 均与 $E$ 相交，因为 $F(0)=0<1\leqslant F(1)$
  \item[\textbf{23.}] (1) $[3,6]$。由 $\dfrac13x\leqslant9\leqslant3x$ 得 $3\leqslant x\leqslant27$，与 $\dfrac23\leqslant x\leqslant6$ 取交；(2) $\left[\dfrac13,2\right]$。数列条件给出 $\dfrac13\leqslant q\leqslant3$，和的条件化为 $q^n\leqslant2S_n$，分 $q=1$、$q>1$、$q<1$ 讨论；(3) $k$ 的最大值为 $1999$，此时公差 $d=-\dfrac1{1999}$。$d\geqslant0$ 时 $k\leqslant1000$；$d<0$ 时由 $d\geqslant-\dfrac2{2k-1}$ 得 $k^2-2000k+1000\leqslant0$
\end{itemize}
